% Fractus White Paper v2.0 — arXiv LaTeX source
% Compile: pdflatex main.tex && bibtex main && pdflatex main.tex && pdflatex main.tex
%
% To submit to arXiv:
% 1. Create account at arxiv.org
% 2. Upload main.tex + any figures
% 3. Select: cs.AI (Artificial Intelligence) + cs.LG (Machine Learning)
% 4. Title: "Fractus: A Continuous Thought Engine with Multi-Block Depth, Self-Modification, and Progressive Growth"
% 5. Abstract: copy from the \begin{abstract} below
% 6. Authors: Philippe-Antoine Robert (rpa.tu@proton.me)

\documentclass[11pt,a4paper]{article}
\usepackage[utf8]{inputenc}
\usepackage[margin=1in]{geometry}
\usepackage{amsmath,amssymb}
\usepackage{booktabs}
\usepackage{graphicx}
\usepackage{hyperref}
\usepackage{url}
\usepackage{array}
\usepackage{xcolor}
\usepackage{listings}

\hypersetup{
    colorlinks=true,
    linkcolor=blue!70!black,
    urlcolor=blue!70!black,
    citecolor=blue!70!black,
}

\title{\textbf{Fractus: A Continuous Thought Engine with Multi-Block Depth, \\
       Self-Modification, and Progressive Growth}}
\author{Philippe-Antoine Robert\\
        \texttt{rpa.tu@proton.me}}
\date{August 6, 2026}

\begin{document}
\maketitle

\begin{abstract}
I present Fractus v2.0 --- a continuous cognitive agent architecture that departs fundamentally from the transformer paradigm. Unlike static models that map input to output in a single forward pass, Fractus is a dynamical system that maintains a persistent thought state, advances it tick by tick through a multi-block residual stack, and emits output only when it has something confident to say. This version introduces three structural advances: (1) multi-block depth with per-block continuous attention state, (2) progressive growth via zero-padding inheritance, and (3) runtime self-modification. Training optimizations (sparse low-rank MoE, head-partial training, gradient accumulation) achieve 707 tokens/second on a consumer CPU. I also report negative results: Expert Decoupled Training and the Forward-Forward algorithm were both tested and refuted. All code is open-source (MIT).
\end{abstract}

\section{Introduction}
Contemporary LLMs are static, stateless, generic, and centralized. Fractus challenges each: the Continuous Thought Engine (CTE) is dynamical, the Persistent Memory gives cross-session recall, the PhaseRoutedMoE enables specialization, and progressive growth enables consumer-hardware training.

\section{Architecture}

\subsection{The Continuous Thought Engine}
The CTE stacks $N$ blocks. Each block owns its attention state $(S, z)$, Kuramoto phases, and PhaseRoutedMoE. The thought $h$ flows through the stack as a residual stream.

\subsection{PhaseRoutedMoE}
Routing via von Mises gate on Farey-distributed expert phases. Sparse gather-first dispatch (top-$k$). Low-rank experts: $W = \text{scale} \cdot U V^\top$ (rank 64).

\subsection{Progressive Growth}
Grow palier by palier (width + depth + experts). Zero-pad old weights. Warm start converges faster.

\section{Training}
Online training: 32 tokens/chunk, 1 backward/chunk, detached state carry. Combined optimizations: tied head, head-partial, sparse MoE, gradient accumulation, Kuramoto detachment. Measured: 707 tok/s on CPU.

\section{Negative Results}
\subsection{EDT --- Refuted}
5 variants tested on 13M CTE. All $\sim$19\% worse than from-scratch. Root cause: objective misalignment + routing concentration.
\subsection{Forward-Forward --- Refuted}
Goodness signal $\neq$ CE. NLL increased from 124 to 221.

\section{Results}
Progressive growth: 4 paliers (6.6M $\to$ 350M) on CPU in 4h44m. Self-modification validated (50\% traffic, stable loss). Cross-session memory verified.

\section{Conclusion}
Fractus v2.0 demonstrates that a continuous, multi-block, self-modifying cognitive agent can be trained on consumer hardware. The architecture is validated by 28 tests and measured benchmarks.

\bibliographystyle{plain}
\begin{thebibliography}{8}
\bibitem{katharopoulos2020} Katharopoulos et al. (2020). Transformers are RNNs. ICML.
\bibitem{siren} Sitzmann et al. (2020). Implicit Neural Representations with SIREN. NeurIPS.
\bibitem{hinton2022} Hinton, G. (2022). The Forward-Forward Algorithm.
\bibitem{kuramoto} Kuramoto, Y. (1984). Chemical Oscillations, Waves, and Turbulence.
\bibitem{lora} Hu et al. (2021). LoRA: Low-Rank Adaptation. arXiv.
\end{thebibliography}

\end{document}
